% --- Archon physics formalization source begin ---
% archon:physics
% archon:chemistry
% archon:covers IChO2026Problems/problem_icho_2026_t4_a8.lean
% archon:source-report reports/icho_2026/problem_icho_2026_t4_a8.source.json
% archon:problem-id icho_2026_t4
% archon:part-id T4-A8
% archon:previous-part icho_2026_t4_a7
% archon:previous-part-policy natural_language_prerequisite_only

\chapter{Icho Chemistry problem icho\_2026\_t4\_a8}
\label{ch:IChO2026Problems_problem_icho_2026_t4_a8}

\paragraph{Problem source.}
\#\# Source
58th International Chemistry Olympiad, Tashkent, Uzbekistan, 2026.
Official-materials index: https://www.icho2026.uz/problems
English problem PDF: ../icho\_2026\_source/raw/theory\_problem.pdf
English solution/rubric PDF: ../icho\_2026\_source/raw/theory\_solution.pdf
Problem PDF page 39; printed local question page 3; solution starts on PDF page 36.
The page PNG is primary visual evidence for formulas, structures, tables, and figures.

\#\# T4. The Nuclear Past of Uzbekistan
T4. The nuclear past of Uzbekistan 6\% Natural uranium consists primarily of two isotopes, 235U and 238U . 235U is responsible for sustaining nuclear fission reactions. 4.1 Calculate the atomic abundance of 235U in natural uranium, assuming it consists only of isotopes 235U (235.04 a.u.) and 238U (238.05 a.u.). The main reaction occurring in nuclear power plants is the reaction where 235U absorbs a neutron and undergoes fission, releasing energy and producing three additional neutrons, thus sustaining a chain reaction. 235 92 U +1 0 n →… + 31 0n Neutrons that induce the chain reaction are produced by several ways: a) Reaction of 11B with 𝛼-particles inside the reactor core b) Interaction of 𝛾rays with 2D nuclei present in the reactor coolant water c) Reaction of 9Be with 𝛼-particles produced in an externally installed neutron source. 4.2 Write the nuclear reaction equations for (a), (b), and (c). The graph above displays the percentage yield, P(A), of fission products by mass number, A, resulting from 235U fission. 4.3 Write the most common nuclear fission equation for splitting of 235U upon neutron absorption. The pair of elements formed in this reaction are in the same group of Periodic Table. 4.4 Calculate the energy released in this reaction (ΔE, MeV), if the binding energy 𝐵𝐸(235U) = 7.59 MeV/nucleon and average binding energy for fission products 𝐵𝐸(fis.) = 8.45 MeV/nucleon. Neglect the binding energy of free neutrons. Fission neutrons are produced with an average energy of 2 MeV and slow down through many scatterings with nuclei. After several collisions, their speed decreases until their average kinetic energy matches that of the surrounding atoms or molecules. The logarithmic energy decrement, 𝜉, shows the average decrease per collision in the logarithm of the neutron energy: 𝜉= ln ( 𝐸𝑖𝑛𝑖𝑡𝑖𝑎𝑙 𝐸𝑓𝑖𝑛𝑎𝑙) 𝐸𝑖𝑛𝑖𝑡𝑖𝑎𝑙and 𝐸𝑓𝑖𝑛𝑎𝑙are the initial and final neutron energies per collision. 𝜉is a constant for each type of material and does not depend on the initial neutron energy. 4.5 Determine the average number of collisions required, n𝑐, to slow a neutron from 2 MeV to 0.012 eV, using water as the moderator. 𝜉(water)=0.948 In 1966, the Urtabulak gas field suffered an uncontrolled natural gas well blowout, releasing enormous volumes of burning methane. 4.6 Calculate Δ𝑟𝐻∘ 298 (kJ mol−1) for one mole of methane combustion reaction at 298 K. Assume all species are gaseous. Use the thermodynamic data below Δ𝑓𝐻∘ 298(CH4) -74.8 kJ mol−1 Δ𝑓𝐻∘ 298(H2O, gas) -241.8 kJ mol−1 Δ𝑓𝐻∘ 298(CO2) -393.5 kJ mol−1 𝐶𝑃(CH4) 35 J mol−1 K−1 𝐶𝑃(H2O, gas) 34 J mol−1 K−1 𝐶𝑃(O2) 29 J mol−1 K−1 𝐶𝑃(CO2) 37 J mol−1 K−1 If you did not get an answer for 4.6, use Δ𝐻298 = −750 𝑘𝐽𝑚𝑜𝑙−1 for further calculations. 4.7 Calculate Δ𝑟𝐻2000 (kJ mol−1) per mole of methane combustion reaction at 2000 K. Assume all species are gaseous. If you did not get an answer for 4.7, use Δ𝐻2000 = −700 𝑘𝐽𝑚𝑜𝑙−1 for further calculations. Assume that methane is flowing out the well at 101.325 kPa and 298 K with volumetric flow Q = 2.2 × 105 m3 before it catches fire.

\#\# Current subquestion 4.8
Calculate the total energy released per day, 𝐸, (in J day−1) by complete isothermic combustion of methane at T = 2000 K

\#\# Dataset metadata
IChO 2026; T4; raw rubric points=2; kind=theory; formalization\_ready=true.

\paragraph{Shared problem context.}
T4. The nuclear past of Uzbekistan 6\% Natural uranium consists primarily of two isotopes, 235U and 238U . 235U is responsible for sustaining nuclear fission reactions. 4.1 Calculate the atomic abundance of 235U in natural uranium, assuming it consists only of isotopes 235U (235.04 a.u.) and 238U (238.05 a.u.). The main reaction occurring in nuclear power plants is the reaction where 235U absorbs a neutron and undergoes fission, releasing energy and producing three additional neutrons, thus sustaining a chain reaction. 235 92 U +1 0 n →… + 31 0n Neutrons that induce the chain reaction are produced by several ways: a) Reaction of 11B with 𝛼-particles inside the reactor core b) Interaction of 𝛾rays with 2D nuclei present in the reactor coolant water c) Reaction of 9Be with 𝛼-particles produced in an externally installed neutron source. 4.2 Write the nuclear reaction equations for (a), (b), and (c). The graph above displays the percentage yield, P(A), of fission products by mass number, A, resulting from 235U fission. 4.3 Write the most common nuclear fission equation for splitting of 235U upon neutron absorption. The pair of elements formed in this reaction are in the same group of Periodic Table. 4.4 Calculate the energy released in this reaction (ΔE, MeV), if the binding energy 𝐵𝐸(235U) = 7.59 MeV/nucleon and average binding energy for fission products 𝐵𝐸(fis.) = 8.45 MeV/nucleon. Neglect the binding energy of free neutrons. Fission neutrons are produced with an average energy of 2 MeV and slow down through many scatterings with nuclei. After several collisions, their speed decreases until their average kinetic energy matches that of the surrounding atoms or molecules. The logarithmic energy decrement, 𝜉, shows the average decrease per collision in the logarithm of the neutron energy: 𝜉= ln ( 𝐸𝑖𝑛𝑖𝑡𝑖𝑎𝑙 𝐸𝑓𝑖𝑛𝑎𝑙) 𝐸𝑖𝑛𝑖𝑡𝑖𝑎𝑙and 𝐸𝑓𝑖𝑛𝑎𝑙are the initial and final neutron energies per collision. 𝜉is a constant for each type of material and does not depend on the initial neutron energy. 4.5 Determine the average number of collisions required, n𝑐, to slow a neutron from 2 MeV to 0.012 eV, using water as the moderator. 𝜉(water)=0.948 In 1966, the Urtabulak gas field suffered an uncontrolled natural gas well blowout, releasing enormous volumes of burning methane. 4.6 Calculate Δ𝑟𝐻∘ 298 (kJ mol−1) for one mole of methane combustion reaction at 298 K. Assume all species are gaseous. Use the thermodynamic data below Δ𝑓𝐻∘ 298(CH4) -74.8 kJ mol−1 Δ𝑓𝐻∘ 298(H2O, gas) -241.8 kJ mol−1 Δ𝑓𝐻∘ 298(CO2) -393.5 kJ mol−1 𝐶𝑃(CH4) 35 J mol−1 K−1 𝐶𝑃(H2O, gas) 34 J mol−1 K−1 𝐶𝑃(O2) 29 J mol−1 K−1 𝐶𝑃(CO2) 37 J mol−1 K−1 If you did not get an answer for 4.6, use Δ𝐻298 = −750 𝑘𝐽𝑚𝑜𝑙−1 for further calculations. 4.7 Calculate Δ𝑟𝐻2000 (kJ mol−1) per mole of methane combustion reaction at 2000 K. Assume all species are gaseous. If you did not get an answer for 4.7, use Δ𝐻2000 = −700 𝑘𝐽𝑚𝑜𝑙−1 for further calculations. Assume that methane is flowing out the well at 101.325 kPa and 298 K with volumetric flow Q = 2.2 × 105 m3 before it catches fire.

\paragraph{Current subquestion.}
Calculate the total energy released per day, 𝐸, (in J day−1) by complete isothermic combustion of methane at T = 2000 K

\paragraph{Recorded answer/context.}
Using the flow as a gas volume measured at 1 atm and 298 K: 𝑛= 𝑃𝑉 𝑅𝑇= 101325 Pa×2.2×105 m3 8.314 J mol−1 K−1×298 K = 9.00 × 106 mol day−1 Total combustion energy released daily = 7.04 × 1012 J day−1 For provided Δ𝑟𝐻2000 = −700 kJ mol−1 the answer = 6.30 × 1012 J day−1 2 points for the correct answer. 1 point if the answer is incorrect but the relevant calculation exists. 2.0 pt

\paragraph{Figure/image paths.}
\begin{itemize}
\item ../icho\_2026\_source/image/T4\_page-3.png
\item ../icho\_2026\_source/image/T4\_page-2.png
\end{itemize}

\paragraph{Reusable previous-part conclusions.}
\begin{itemize}
\item Source T4-A7. Question: Calculate Δ𝑟𝐻2000 (kJ mol−1) per mole of methane combustion reaction at 2000 K. Assume all species are gaseous. Reusable conclusions: Assume constant heat capacities. Δ𝐻(𝑇) = Δ𝐻(298 𝐾) + Δ𝐶𝑝(𝑇−298 𝐾) Δ𝐶𝑝= [𝐶𝑝(CO2) + 2𝐶𝑝(H2O)] −[𝐶𝑝(CH4) + 2𝐶𝑝(O2)] Δ𝐶𝑝= [37 + 2(34)] −[35 + 2(29)] = 12 J mol−1 K−1 Δ𝑟𝐻∘ 2000 = −802.3 kJ mol−1 + (12Jmol−1K−1)(1702 𝐾) 1000 = −781.876 kJ mol−1 Δ𝑟𝐻∘ 2000 = −781.9 kJ mol−1 For a given value of Δ𝑟𝐻∘ 298 = −750 kJ/mol alternative answer will be: Δ𝑟𝐻2000 = −781.9 𝑘𝐽𝑚𝑜𝑙−1 Δ𝑟𝐻∘ 2000 = −750 kJ mol−1 + (12 J mol−1K−1)(1702 K) 1000 = −729.576 kJ mol−1 Δ𝑟𝐻2000 = −729.6 kJ mol−1 2 points total: 1 point for the correct heat capacity change calculation. 1 point for the correct enthalpy calculation at 2000 K. 2.0 pt Policy: natural\_language\_prerequisite\_only; the current target and later answers must not be assumed
\end{itemize}

\paragraph{Formalization target.}
create a compiling Lean file with sorry bodies at `IChO2026Problems/problem\_icho\_2026\_t4\_a8.lean`.
The Lean declarations must preserve every mathematical object, quantifier, hypothesis, side condition, approximation/error guarantee, and requested conclusion from the source statement.
Use Mathlib/Physlib/Chemistry names found through LeanExplore where available. Prefer the configured benchmark Base library; introduce a local definition only when the source genuinely requires an object that the environment does not expose.

\begin{theorem}[Icho Chemistry formalization target]
\label{thm:physics:icho_2026_t4_a8:target}
The assigned autoformalize agent should translate this IChO chemistry problem into Lean declarations in the covered file, with theorem and lemma proof bodies written as `by sorry`.
\end{theorem}
\begin{proof}
This is an autoformalization task, not a proof task. Produce faithful statements that can later be proved without weakening the source contract.
\end{proof}
% --- Archon physics formalization source end ---
